\section{实验设计：把“有效”变成可复现证据}

\subsection{实验矩阵设计}
课程建议将 Agent 实验拆成 ``能力变量'' 与 ``系统变量'' 两组：能力变量包括 reasoning budget、memory 策略、planning 粒度；系统变量包括工具延迟、缓存命中率、并发策略和失败重试规则。只有同时控制两组变量，实验结论才有解释力。

\begin{knowledgebox}{最小可复现实验单元}
\begin{itemize}
    \item 固定任务集合与数据版本；
    \item 固定工具接口与权限边界；
    \item 固定模型版本与温度参数；
    \item 记录完整轨迹与中间状态。
\end{itemize}
\end{knowledgebox}

\begin{table}[H]
\centering
\begin{tabular}{p{3.1cm}p{3.3cm}p{2.7cm}p{2.2cm}}
\toprule
\textbf{实验项} & \textbf{对照设置} & \textbf{主指标} & \textbf{副指标} \\
\midrule
Reasoning budget & 2k / 6k / 12k tokens & 任务成功率 & 单任务成本 \\
Memory policy & 无记忆 / 检索记忆 / 分层记忆 & 长任务完成率 & 错误恢复时长 \\
Planning policy & 静态计划 / 动态重规划 & 平均步数 & 无效动作率 \\
Verification depth & 终局校验 / 步骤校验 / 双层校验 & 正确率 & 延迟 \\
\bottomrule
\end{tabular}
\caption{Agent 消融实验的基础矩阵}
\end{table}

\subsection{失败分类与归因方法}
没有统一失败分类，实验就会陷入 ``调参玄学''。推荐把失败划分为四类：感知失败、推理失败、执行失败、验证失败。每类失败对应不同修复动作，避免 ``所有问题都加大模型'' 的低效路径。

\begin{importantbox}{失败归因要做到“可执行”}
归因结果应能直接映射到工程动作：
\begin{itemize}
    \item 感知失败 $\rightarrow$ 强化解析器与页面状态检查；
    \item 推理失败 $\rightarrow$ 调整 reasoning 模板与反思触发；
    \item 执行失败 $\rightarrow$ 工具封装与重试策略优化；
    \item 验证失败 $\rightarrow$ 扩展校验器覆盖与错误阈值。
\end{itemize}
\end{importantbox}

\begin{warningbox}{避免“成功样本偏见”}
如果只分析成功轨迹，系统会误判自身鲁棒性。应优先分析高损失失败样本（高成本、长链路、不可恢复），这些样本决定生产可用性上限。
\end{warningbox}

\subsection{线上迭代节奏}
课程并不建议一次性上线大改。更稳妥的路径是 ``影子模式（shadow mode）$\rightarrow$ 小流量灰度 $\rightarrow$ 全量发布''。每个阶段都应有明确回滚条件，并持续监控任务级别 SLO。

\begin{knowledgebox}{线上监控建议}
\begin{itemize}
    \item 成功率下探阈值与自动降级策略；
    \item 成本突增阈值与预算闸门；
    \item 工具错误率阈值与熔断机制。
\end{itemize}
\end{knowledgebox}

\subsection{本章小结}
实验设计的目标不是“证明系统很好”，而是“定位系统为何好/为何坏”。可复现、可归因、可回滚是 Agent 实验体系的三条底线。

